\subsection{Comparación con el perfil continuo: inversa uniforme y conclusión}\label{sec:P06}
Se usan (TR1)--(TR9). Las constantes dependen de las cotas de la Sección~\ref{sec:P02} y
del perfil continuo, pero no de $n$.

\paragraph{Jacobiana media exacta.}
Derivar la raíz de (TR5) da, para cada casilla,
\[
 w(z)_{ij}=\frac{\partial\rho(z)_{ij}}{\partial(\mathcal Tz)_{ij}}
 =\frac{\rho(z)_{ij}(1-\rho(z)_{ij}/n)}{1-4\rho(z)_{ij}/n}.
\]
En el segmento de (TR9), $a\le w(z)\le C_w:=2Q$. La jacobiana es
\[
 J(z)h=\operatorname{diag}(0_R,0_C,u_D(z),u_A(z))h
                         +\mathcal R(w(z)\mathcal T h).
                                                        \tag{TP1}
\]
Su promedio en $z(t)=(1-t)z^0+tz^*$ conserva exactamente esa forma:
\[
 \overline J=\int_0^1J(z(t))\,dt
 =\operatorname{diag}(0_R,0_C,\overline u_D,\overline u_A)
                              +\mathcal R(\overline w\mathcal T),
 \quad a\le\overline w\le C_w,\quad\zeta\le\overline u_D,\overline u_A\le1.
                                                        \tag{TP2}
\]
Es autoadjunta para los productos escalares definidos entre (TR3) y (TR4), anula $v_0$ y
preserva $E_n$.

\paragraph{Coercividad en el gauge.}
La estimación se prueba directamente para el operador (TP2), que no es la matriz evaluada en el
maximizador; el argumento es el de la Sección~\ref{sec:P03}.
Si $h\in E_n$, las medias de $h_R$ y $h_C$ son iguales. Por ello
\[
 \|h_R\circ R+h_C\circ C\|_\square^2
 =\|h_R\|_{2,R}^2+\|h_C\|_{2,C}^2+2\overline h_R\overline h_C
 \ge\|h_R\|_{2,R}^2+\|h_C\|_{2,C}^2.
\]
Cada levantamiento diagonal tiene norma $L^2$ a lo sumo uno, pues
una diagonal contiene a lo sumo $n$ casillas. Escribiendo
$h_R\circ R+h_C\circ C=\mathcal T h-h_D\circ D-h_A\circ A$
y usando $(s-d-a)^2\le2s^2+4(d^2+a^2)$, resulta
\[
 \|h\|_\ell^2\le2\|\mathcal T h\|_\square^2
                    +5(\|h_D\|_{2,D}^2+\|h_A\|_{2,A}^2).
\]
Como $\mathcal R^*=\mathcal T$, (TP2) implica
\[
 \langle h,\overline Jh\rangle_\ell
 \ge a\|\mathcal T h\|_\square^2
             +\zeta(\|h_D\|_{2,D}^2+\|h_A\|_{2,A}^2)
 \ge\lambda\|h\|_\ell^2,
 \quad\lambda=\min(a/2,\zeta/5)>0.
                                                        \tag{TP3}
\]
Por dimensión finita, $\overline J:E_n\to E_n$ es invertible y
$\|\overline J^{-1}k\|_\ell\le\lambda^{-1}\|k\|_\ell$.

\paragraph{De $L^2$ al supremo mediante transversalidad discreta.}
Sea $\overline Jh=k\in E_n$. La ecuación de una línea $\ell$ de familia
$f$ es
\[
 \left(d_{f,\ell}+\frac1n\sum_{ij\in\ell}\overline w_{ij}\right)h_{f,\ell}
 +\sum_{g\ne f}\frac1n\sum_{ij\in\ell}\overline w_{ij}h_{g,\ell_g(i,j)}
 =k_{f,\ell},
                                                        \tag{TP4}
\]
con $d_R=d_C=0$, $d_D=\overline u_D$, $d_A=\overline u_A$.
El coeficiente propio es al menos $a_0:=\min(a,\zeta)>0$:
las filas y columnas tienen exactamente $n$ casillas; las diagonales y antidiagonales tienen
$d_{f,\ell}\ge\zeta$.
Dos líneas de familias distintas se cortan en a lo sumo una casilla.
Así, para $g\ne f$, los índices $\ell_g(i,j)$ visitados sobre $\ell$
no se repiten. Cauchy--Schwarz da
\[
 \frac1n\sum_{ij\in\ell}|h_{g,\ell_g(i,j)}|
 \le\sqrt{\frac{|\ell|}{n}}
       \left(\frac1n\sum_{\ell'\text{ visitada}}|h_{g,\ell'}|^2\right)^{1/2}
 \le\|h_g\|_{2,g}.
                                                        \tag{TP5}
\]
En el acoplamiento entre diagonales y antidiagonales se visita sólo una paridad, pero sigue siendo
un subconjunto sin repeticiones: la misma cota vale sin factor adicional.
Por (TP4)--(TP5) y la suma de las tres familias transversales,
\[
 \|h\|_\infty\le a_0^{-1}
       (\|k\|_\infty+\sqrt3 C_w\|h\|_\ell).
\]
Como hay $6n-2$ líneas, $\|k\|_\ell\le\sqrt6\|k\|_\infty$.
Con (TP3) se obtiene la inversa uniforme que necesitamos:
\[
 \|\overline J^{-1}k\|_\infty\le K\|k\|_\infty,
 \qquad K=a_0^{-1}(1+\sqrt{18}C_w/\lambda).
                                                        \tag{TP6}
\]
La estimación admite datos arbitrarios en $E_n$, incluso oscilatorios.

\begin{proposicion}
Para $n$ suficientemente grande, el maximizador canónico satisface
\[
 \max_{ij}|n\pi_{ij}-\nu_1(x_i,x_j)|\le C/n,
 \quad C=4C_wKC_F+12C_\nu^2,
                                                        \tag{TP7}
\]
\[
 \max_k|h^*_{D,k}-u_D(k/n)|\le KC_F/n,\qquad
 \max_l|h^*_{A,l}-u_A(l/n)|\le KC_F/n.
                                                        \tag{TP8}
\]
\end{proposicion}
\begin{proof}
El teorema fundamental del cálculo y (TR5) dan, exactamente,
$-F_n(z^0)=F_n(z^*)-F_n(z^0)=\overline J(z^*-z^0)$.
Ambos potenciales están en $E_n$; (TR8) y (TP6) implican
$\|z^*-z^0\|_\infty\le KC_F/n$.
En el segmento, $d\rho/d(\mathcal Tz)=w\le C_w$ y
$\|\mathcal T\|_{\infty\to\infty}\le4$. Por tanto
$\|\rho(z^*)-\rho(z^0)\|_\infty\le4C_wKC_F/n$.
(TR7) y la identidad $\rho(z^*)=n\pi$ prueban (TP7).
Las vacancias exponenciales son a lo sumo uno a lo largo del segmento;
el teorema del valor medio escalar da (TP8).
\end{proof}

Así, (TP7) establece (LC3) con el mismo muestreo de puntos medios.
